\section{Related Work}
\label{sec:related_work}

\subsection{Continual learning methods}
\label{sec:rw_cl}

\textbf{Regularization-based methods} reduce forgetting by penalizing changes to parameters considered important for earlier tasks, without allocating separate modules per task. EWC \citep{kirkpatrick2017overcoming} estimates parameter importance from the Fisher information matrix and penalizes deviation from previous values; Learning without Forgetting \citep{li2017learning} instead preserves the model's output distribution via knowledge distillation; Synaptic Intelligence \citep{zenke2017continual} and Memory Aware Synapses \citep{aljundi2018memory} accumulate importance online. These methods are memory-efficient because they require no replay buffer, but because all tasks share one parameter space, their protection weakens as the number and diversity of tasks grows, which Table~\ref{tab:codetask_qwen} confirms for EWC on the eight-task CodeTask sequence. GMM-Router sidesteps this failure mode by allocating a dedicated LoRA adapter per task (Section~\ref{sec:experts}) rather than sharing one regularized parameter set.

\textbf{Memory-based methods} store information from earlier tasks and replay it during later training. GEM and A-GEM \citep{lopezpaz2017gem,chaudhry2019agem} use stored examples to constrain gradients so earlier tasks are not harmed; DER++ \citep{buzzega2020dark} replays stored logits to match past model behavior; Deep Generative Replay \citep{shin2017continual} generates pseudo-examples instead of storing real ones. These are strong baselines but increase per-step cost and conflict with settings, such as proprietary or privacy-sensitive code repositories, where retaining historical data is undesirable (Section~\ref{sec:motivation}). GMM-Router requires no stored training examples at inference time: its GMMs are fitted once from task-level routing embeddings during training and then frozen.

\textbf{Architecture-based methods} handle forgetting by partitioning model capacity across tasks. Progressive Neural Networks add a branch per task, PackNet reuses a fixed network through pruning, and HAT learns hard attention masks \citep{rusu2016progressive,mallya2018packnet,serra2018overcoming}; L2P \citep{wang2022learning} learns a shared, dynamically retrieved prompt pool, and Progressive Prompts \citep{razdaibiedina2023progressive} keeps earlier prompts fixed while adding new ones per task. These are included as baselines in Section~\ref{sec:baselines} (L2P, ProgPrompt); their limited representational capacity relative to full LoRA experts is consistent with their weaker CodeTask and executable-benchmark results (Tables~\ref{tab:codetask_qwen}--\ref{tab:exec_avg}).

\textbf{LoRA-based continual learning} is the branch most directly related to GMM-Router. A single LoRA trained sequentially (SeqLoRA in this paper) keeps storage flat but allows forgetting because all tasks share the same low-rank update \citep{hu2022lora}; per-task LoRAs (the LoRA-per-task upper bound) prevent forgetting but require storage that grows linearly and either task identity or a router at inference. O-LoRA \citep{wang2023orthogonal} reduces interference by learning task updates in mutually orthogonal low-rank subspaces; GORP \citep{wang2025continual} uses low-rank gradient projection to expand the optimization space; TreeLoRA \citep{qian2025treelora} organizes adapters hierarchically by gradient similarity. CoLoR \citep{wistuba2023continual} studies whether LoRA itself is an effective continual learning mechanism for pretrained transformers; C-LoRA \citep{zhang2025clora} introduces a learnable routing matrix that manages low-rank updates across tasks while encouraging orthogonality; CL-LoRA \citep{he2025cllora} combines task-shared and task-specific adapters with orthogonal structures and distillation for rehearsal-free continual learning. These methods reduce --- but, because they still couple tasks through a partially shared subspace, do not eliminate --- cross-task interference, and the number of usable orthogonal directions shrinks as more tasks are added (Section~\ref{sec:introduction}). GMM-Router instead removes coupling between adapters entirely by isolating each task in its own LoRA module and coordinating them only at inference, through routing rather than through a shared or orthogonalized subspace.

Online-CL-LLM \citep{cheng2025exploiting}, evaluated as a baseline in Section~\ref{sec:baselines}, is the method closest in spirit to GMM-Router: it characterizes each PEFT block by a representative feature distribution and dynamically selects blocks per input without task labels, sharing GMM-Router's reliance on distributional information for task-agnostic routing. It differs in that it selects among existing PEFT blocks rather than combining them, whereas GMM-Router uses explicit per-task GMM posterior scores (Eq.~\eqref{eq:snorm}) to support both hard selection and soft, weight-level adapter merging (Eq.~\eqref{eq:weight_merge}), which Section~\ref{sec:exec_results} shows is responsible for GMM-Router's Pass@5 lead on the executable benchmark.

\subsection{Mixture of experts and density-based routing}
\label{sec:rw_moe}

The routing mechanism in Section~\ref{sec:methodology} is related to the mixture-of-experts literature, in which expert networks specialize in different parts of the input space and a gating function determines their contribution. \citet{jacobs1991adaptive} introduced adaptive mixtures of local experts with a learned gating network, and \citet{jordan1994hierarchical} extended this to hierarchical mixtures trained with EM-style learning, establishing that routing and expert specialization admit a principled probabilistic interpretation. In modern deep learning, sparsely-gated mixture-of-experts layers \citep{shazeer2017outrageously} and Switch Transformers \citep{fedus2022switch} activate only a small number of experts per input using learned neural gating rather than an explicit density model, but the underlying motivation --- routing different inputs to specialized modules --- is the same. GMM-Router differs from this line of work in two respects relevant to the continual setting: its gate is an explicit, per-task Gaussian Mixture Model fitted independently for each task as it arrives (Section~\ref{sec:router}), rather than a single jointly trained gating network that would need to be retrained as new experts are added; and its routing features are extracted from a frozen backbone through a fixed orthonormal projection (Section~\ref{sec:features}) specifically so that GMM densities fitted at different points in the task sequence remain comparable, a stability requirement that does not arise in the single-training-run setting where standard MoE gating is typically deployed.

\subsection{Code language models and evaluation}
\label{sec:rw_code}

Pretrained code language models, reviewed in Section~\ref{sec:clm_background}, provide the backbones on which all methods compared in this paper are built \citep{feng2020codebert,guo2021graphcodebert,guo2022unixcoder,ahmad2021plbart,wang2021codet5,wang2023codet5plus,fried2023incoder,roziere2023codellama,lozhkov2024starcoder2,hui2024qwen25coder}. Continual learning for this domain was introduced by CodeTask-CL \citep{yadav2023exploring}, which this paper extends from four to eight tasks (Section~\ref{sec:codetask_data}) because a four-task sequence scored only by BLEU \citep{papineni2002bleu} and CodeBLEU \citep{ren2020codebleu} under-represents both the length of realistic deployment sequences and the risk that text-similarity metrics reward superficially correct but functionally broken code. Execution-based evaluation, following the Codex/HumanEval methodology of \citet{chen2021evaluating} and the multilingual MultiPL-E benchmark \citep{cassano2023multipl}, addresses the second gap directly but has not previously been combined with sequential, task-agnostic continual learning; the executable benchmark and deduplication pipeline in Section~\ref{sec:data} are constructed specifically to close this gap.
